\section*{Chapter Summary}

This chapter established the methodological foundations of the thesis.
The Double Diamond model was introduced as the overarching process framework, mapping the four phases of the research (literature review, requirements analysis, design and implementation, and evaluation) onto two sequential diverge-then-converge cycles covering the problem space and the solution space respectively.
The research approach was characterised as artifact-oriented, treating the constructed artifact and its architectural decisions as primary research outputs.
The engineering process was described as iterative and vertical-slice driven, with continuous integration, pull-request-based review, and open-source publication supporting reproducibility and extensibility.
Requirements were shown to be derived through retrospective analysis of an existing prototype, refined through supervision sessions, and expressed as user stories and MoSCoW-prioritised requirements.
The evaluation strategy was defined around two questions: functional correctness verified against the requirements, and architectural quality assessed through dependency inspection and external contract validation; a structural threat to validity was acknowledged and mitigated through open publication of the codebase.
Ethical responsibilities arising from the collection of gaze and participant data were identified and assigned to the researcher operating the platform.
The requirements derived through the process described here are presented in full in the following chapter.
